\subsection{CAN: একই Campus-এর LAN-গুলো যুক্ত করা}
\begin{definitionbox}{Campus Area Network (CAN)}
একই শিক্ষাপ্রতিষ্ঠান বা প্রতিষ্ঠানের ক্যাম্পাসের একাধিক ভবনের LAN সংযুক্ত করে গঠিত নেটওয়ার্ককে Campus Area Network বলে। যেমন, কলেজের ল্যাব, গ্রন্থাগার ও প্রশাসনিক ভবনের LAN ফাইবার ব্যাকবোন দিয়ে যুক্ত করা।
\end{definitionbox}
PAN ব্যক্তির কাছের ডিভাইস, LAN স্থানীয় কক্ষ বা ভবন, CAN একই ক্যাম্পাসের কয়েকটি LAN, MAN একটি মহানগর এলাকা এবং WAN বৃহৎ ভৌগোলিক এলাকার নেটওয়ার্ক বোঝায়। একটি নির্দিষ্ট দূরত্ব দিয়ে সব শ্রেণির সীমা অপরিবর্তনীয়ভাবে ভাগ করা যায় না। শ্রেণি নির্ধারণে বিস্তৃতি বিবেচ্য; শুধু ডেটা হার যথেষ্ট নয়। Ethernet ও Wi-Fi প্রযুক্তির নাম, Star সংযোগের বিন্যাস এবং LAN ভৌগোলিক বিস্তৃতির শ্রেণি। একই নেটওয়ার্কে এই তিন পরিচয় প্রযোজ্য হতে পারে।
\textbf{আরও ব্যবহারভিত্তিক নাম:} HAN (Home Area Network) বাসার কম্পিউটার, প্রিন্টার ও স্মার্ট ডিভাইসের নেটওয়ার্ক; এর বিস্তৃতি সাধারণত স্থানীয়। SAN (Storage Area Network) সার্ভারকে সংরক্ষণসম্পদের সঙ্গে যুক্ত করার বিশেষ নেটওয়ার্ক। যেমন, ডেটা সেন্টারে সার্ভার ও ডিস্ক স্টোরেজের সংযোগ। এখানে Home ও Storage পরিবেশ এবং উদ্দেশ্য বোঝায়। এগুলো PAN/LAN/MAN/WAN-এর দূরত্বভিত্তিক শ্রেণির নতুন ধাপ নয়।

\subsection{মালিকানা, দায়িত্ব ও নিয়ন্ত্রণ আলাদা ভিত্তি}
\begin{tabularx}{\linewidth}{|p{.23\linewidth}|X|X|}
\hline
\textbf{ভিত্তি} & \textbf{ধরন} & \textbf{অর্থ ও উদাহরণ}\\\hline
মালিকানা ও প্রবেশাধিকার & Private Network & প্রতিষ্ঠানের নির্ধারিত ব্যবহারকারীর জন্য নিজস্ব নেটওয়ার্ক; যেমন অফিসের LAN বা Intranet।\\\hline
মালিকানা ও প্রবেশাধিকার & Public Network & সাধারণ গ্রাহক সেবাদাতার নেটওয়ার্ক ব্যবহার করতে পারে। সবার ডেটা উন্মুক্ত থাকে—এমন অর্থ নয়।\\\hline
সেবা প্রদানের দায়িত্ব & Client–server & নির্দিষ্ট সার্ভার সেবা দেয়; ক্লায়েন্ট সেবা চায় ও ব্যবহার করে। প্রয়োজনমতো একাধিক সার্ভার থাকতে পারে।\\\hline
সেবা প্রদানের দায়িত্ব & Peer-to-peer & একই পিয়ার প্রয়োজনমতো ক্লায়েন্ট বা সার্ভারের ভূমিকা নেয়। নিবেদিত কেন্দ্রীয় সার্ভার অপরিহার্য নয়।\\\hline
সেবা প্রদানের দায়িত্ব & Hybrid Architecture & কিছু সেবা কেন্দ্রীয় সার্ভার থেকে এবং কিছু পিয়ারগুলোর মধ্যে পাওয়া যায়; যেমন কেন্দ্রীয় লগইন ও পিয়ারদের মধ্যে ফাইল ভাগাভাগি।\\\hline
\end{tabularx}
\textbf{কেন্দ্রীভূত ও বণ্টিত ব্যবস্থা:} তথ্য, সিদ্ধান্ত ও প্রক্রিয়াকরণ প্রধানত একটি কেন্দ্রে হলে তাকে Centralized ব্যবস্থা বলে। একাধিক সংযুক্ত ব্যবস্থায় কাজ ভাগ করে দেওয়া হলে তাকে Distributed ব্যবস্থা বলে। বণ্টিত ব্যবস্থায়ও কেন্দ্রীয় লগইন থাকতে পারে। Client–server-এর ভূমিকা, প্রক্রিয়াকরণের স্থান এবং কেবলের সংযোগবিন্যাস আলাদা ভিত্তি।
\textbf{তিনটি Hybrid:} Hybrid Architecture সেবা প্রদানের মিশ্র দায়িত্ব বোঝায়। Hybrid Topology বিভিন্ন সংযোগবিন্যাসের সমন্বয়। Hybrid Cloud পৃথক ক্লাউড অবকাঠামো সংযুক্ত করে গঠিত সমন্বিত ব্যবস্থা। একই শব্দের অর্থ এখানে শ্রেণিবিন্যাসের ভিত্তি অনুযায়ী আলাদা।

\textbf{নেটওয়ার্ক ব্যবহার মানেই কম খরচ নয়:} প্রিন্টার ভাগাভাগি করলে হার্ডওয়্যার ব্যয় কমে। তবে কেবল, সুইচ, সার্ভার, ব্যাকআপ, নিরাপত্তা ও দক্ষ কর্মীর জন্য বাড়তি খরচ হয়। দ্রুত কাজ, নিয়ন্ত্রণ ও নির্ভরযোগ্যতার সুবিধার সঙ্গে মোট ব্যয় তুলনা করতে হয়।

\subsection{Computer Network-এর বিকাশ: মানুষ, কাজ ও সাল}
\begin{tabularx}{\linewidth}{|p{.15\linewidth}|p{.31\linewidth}|X|}\hline
\textbf{সাল} & \textbf{মানুষ/ব্যবস্থা} & \textbf{অবদান}\\\hline
১৯৬৯ & ARPANET; একাধিক গবেষণা দল & প্রথম হোস্ট থেকে হোস্টে বার্তা পাঠানো; বছরের শেষে চারটি হোস্ট যুক্ত ছিল\\\hline
১৯৭৩ & Robert Metcalfe ও David Boggs; Xerox PARC & Ethernet-এর বিকাশ; LAN প্রযুক্তির গুরুত্বপূর্ণ ধাপ\\\hline
১৯৭৪ & Vint Cerf ও Robert Kahn & পৃথক প্যাকেট নেটওয়ার্ক যুক্ত করার প্রোটোকল-সংক্রান্ত গবেষণাপত্র\\\hline
১৯৮৩ & ARPANET-এর TCP/IP পরিবর্তন & ইন্টারনেটের বিকাশে গুরুত্বপূর্ণ প্রোটোকল পরিবর্তনের ধাপ\\\hline
\end{tabularx}
Cerf ও Kahn-কে ইন্টারনেটের জনক হিসেবে পরিচিত করা হয়। Ethernet, Internet ও Web-এর সৃষ্টি পৃথক ঘটনা। PAN/MAN/WAN ভৌগোলিক বিস্তৃতির শ্রেণি; প্রতিটির জন্য একক আবিষ্কারের সাল নির্ধারণ করা যায় না।
